\section{Présentation du corpus}

Comme indiqué en introduction, cette étude porte sur un corpus composé exclusivement de romans d'Émile Zola. Le choix a été fait de ne pas inclure les nouvelles, les textes journalistiques ni les essais, afin de conserver un ensemble relativement homogène du point de vue générique. Les romans offrent en effet une matière textuelle plus longue et plus développée, ce qui permet d'observer plus finement les thèmes récurrents, les milieux sociaux et les motifs narratifs présents dans l'œuvre de Zola.

L'exclusion des nouvelles répond également à une contrainte méthodologique. Ces textes sont généralement plus courts que les romans et leur intégration aurait pu créer un déséquilibre dans le corpus. Dans le cadre d'une analyse thématique fondée sur des segments textuels, les différences importantes de longueur entre les œuvres peuvent influencer la représentation des thèmes. Limiter le corpus aux romans permet donc de réduire cet effet et de travailler sur un ensemble plus comparable.

Le corpus comprend les œuvres suivantes :

\begin{longtable}{p{1.6cm}p{4.2cm}p{8cm}}
\caption{Corpus des œuvres de Zola utilisées dans l'analyse}\\
\toprule
\textbf{Date} & \textbf{Titre} & \textbf{Fichier} \\
\midrule
\endfirsthead

\toprule
\textbf{Date} & \textbf{Titre} & \textbf{Fichier} \\
\midrule
\endhead

1865 & \textit{La confession de Claude} & \texttt{1865\_La\_confession\_de\_Claude.\_clean.txt} \\
1866 & \textit{Le vœu d'une morte} & \texttt{1866\_Le\_voeu\_d\_une\_morte.\_clean.txt} \\
1867 & \textit{Les mystères de Marseille} & \texttt{1867\_Les\_mysteres\_de\_Marseille.\_clean.txt} \\
1867 & \textit{Thérèse Raquin} & \texttt{1867\_Therese\_Raquin.\_clean.txt} \\
1868 & \textit{Madeleine Férat} & \texttt{1868\_Madeleine\_Ferat.\_clean.txt} \\
1871 & \textit{La fortune des Rougon} & \texttt{1871\_1\_La\_fortune\_des\_Rougon.\_clean.txt} \\
1872 & \textit{La curée} & \texttt{1871\_2\_La\_curee.\_clean.txt} \\
1873 & \textit{Le ventre de Paris} & \texttt{1873\_3\_Le\_ventre\_de\_Paris.\_clean.txt} \\
1874 & \textit{La conquête de Plassans} & \texttt{1874\_4\_La\_conquete\_de\_Plassans.\_clean.txt} \\
1875 & \textit{La faute de l'abbé Mouret} & \texttt{1875\_5\_La\_faute\_de\_l\_abbe\_Mouret.\_clean.txt} \\
1876 & \textit{Son Excellence Eugène Rougon} & \texttt{1876\_6\_Son\_Excellence\_Eugene\_Rougon.\_clean.txt} \\
1877 & \textit{L'assommoir} & \texttt{1877\_7\_L\_assommoir.\_clean.txt} \\
1878 & \textit{Une page d'amour} & \texttt{1878\_8\_Une\_page\_d\_amour.\_clean.txt} \\
1880 & \textit{Nana} & \texttt{1880\_9\_Nana.\_clean.txt} \\
1882 & \textit{Pot-Bouille} & \texttt{1882\_10\_Pot-bouille.\_clean.txt} \\
1883 & \textit{Au Bonheur des dames} & \texttt{1883\_11\_Au\_Bonheur\_des\_dames.\_clean.txt} \\
1884 & \textit{La joie de vivre} & \texttt{1884\_12\_La\_joie\_de\_vivre.\_clean.txt} \\
1885 & \textit{Germinal} & \texttt{1885\_13\_Germinal.\_clean.txt} \\
1886 & \textit{L'œuvre} & \texttt{1886\_14\_L\_oeuvre.\_clean.txt} \\
1887 & \textit{La terre} & \texttt{1887\_15\_La\_terre.\_clean.txt} \\
1888 & \textit{Le rêve} & \texttt{1888\_16\_Le\_reve.\_clean.txt} \\
1890 & \textit{La bête humaine} & \texttt{1890\_17\_La\_bete\_humaine.\_clean.txt} \\
1891 & \textit{L'argent} & \texttt{1891\_18\_L\_argent.\_clean.txt} \\
1892 & \textit{La débâcle} & \texttt{1892\_19\_La\_debacle.\_clean.txt} \\
1893 & \textit{Le docteur Pascal} & \texttt{1893\_20\_Le\_docteur\_Pascal.\_clean.txt} \\
1894 & \textit{Lourdes} & \texttt{1894\_1\_Lourdes.\_clean.txt} \\
1896 & \textit{Rome} & \texttt{1896\_2\_Rome.\_clean.txt} \\
1898 & \textit{Paris} & \texttt{1898\_3\_Paris.\_clean.txt} \\
1899 & \textit{Fécondité} & \texttt{1899\_1\_Fecondite.\_clean.txt} \\
1901 & \textit{Travail} & \texttt{1901\_2\_Travail.\_clean.txt} \\
1903 & \textit{Vérité} & \texttt{1903\_3\_Verite.\_clean.txt} \\

\bottomrule
\end{longtable}

Le corpus est composé de 31 romans, qui couvrent les principales périodes de la production romanesque de Zola. Il comprend :

\begin{itemize}
    \item 5 romans antérieurs au cycle des \textit{Rougon-Macquart} ;
    \item 20 romans appartenant au cycle des \textit{Rougon-Macquart} ;
    \item 3 romans du cycle des \textit{Trois Villes} ;
    \item 3 romans publiés du cycle des \textit{Quatre Évangiles}.
\end{itemize}

Pour ce qui est des \textit{Quatre Évangiles}, \textit{Vérité} paraît à titre posthume. \textit{Justice} ne paraîtra jamais, l'ouvrage étant resté à l'état d'ébauche au moment de la mort de l'écrivain.

Ce corpus permet donc de couvrir une large période de la carrière de Zola, depuis ses premiers romans jusqu'à ses dernières œuvres. Cette amplitude chronologique est importante pour l'analyse, car elle permet d'observer non seulement les thèmes récurrents de l'œuvre romanesque, mais aussi leur possible évolution au fil du temps\footnote{En vue d'une éventuelle analyse ultérieure par dynamic topic modeling sur ce même corpus.}.